\ifdefined\gkver\else\def\gkver{student}\fi
\documentclass[\gkver]{gaokaozhenti}
\begin{document}

\chapter{1993年上海}

\begin{enumerate}
\item
%% number: 1
%% paperName: 1993年上海高考第 1 题 4 分
%% typeId: 1
%% type: 单选题
%% chapter: 原子和原子核
%% point: 人工核反应
%% method: 无
%% score: 4
%% degree: 750
%% duplicateId: 0
%% body:
以下几个原子核反应中，X 代表 $\alpha$ 粒子的反应式是\xzanswer{C}
\fourchoices[answer=C]
{\ce{^{4}_{2}He + ^{9}_{4}Be -> ^{12}_{6}C + X}}
{\ce{^{234}_{90}Th -> ^{234}_{91}Pa + X}}
{\ce{^{2}_{1}H + ^{3}_{1}H -> ^{1}_{0}n + X}}
{\ce{^{30}_{15}P -> ^{30}_{14}Si + X}}
%% answer: C
%% memo:
\memoanswer{
本题原卷及材料中未提供详解，待补充。
}

\item
%% number: 2
%% paperName: 1993年上海高考第 2 题 4 分
%% typeId: 1
%% type: 单选题
%% chapter: 交流电
%% point: 交流电
%% method: 无
%% score: 4
%% degree: 750
%% duplicateId: 0
%% body:
在如图所示电路中，已知交流电源电压 $u = 200\sin 100\pi t\UV$，电阻 $R = 100\UO$，则电流表和电压表的读数分别为\xzanswer{B}
\onepicture{\includesvg{figs/02}}
\fourchoices[answer=B]
{$1.41\UA$，$200\UV$}
{$1.41\UA$，$141\UV$}
{$2\UA$，$200\UV$}
{$2\UA$，$141\UV$}
%% answer: B
%% memo:
\memoanswer{
本题原卷及材料中未提供详解，待补充。
}

\item
%% number: 3
%% paperName: 1993年上海高考第 3 题 4 分
%% typeId: 1
%% type: 单选题
%% chapter: 气体性质
%% point: 气体的状态参量
%% method: 无
%% score: 4
%% degree: 750
%% duplicateId: 0
%% body:
如图所示，U 形管封闭端内有一部分气体被水银封住，已知大气压强为 $p_{0}$，则被封部分气体的压强 $p$（以 $\UcmHg$ 为单位）为\xzanswer{B}
\onepicture{\includesvg{figs/03}}
\fourchoices[answer=B]
{$p_{0} + h_{2}$}
{$p_{0} - h_{1}$}
{$p_{0} - (h_{1} + h_{2})$}
{$p_{0} + h_{2} - h_{1}$}
%% answer: B
%% memo:
\memoanswer{
本题原卷及材料中未提供详解，待补充。
}

\item
%% number: 4
%% paperName: 1993年上海高考第 4 题 4 分
%% typeId: 1
%% type: 单选题
%% chapter: 力物体的平衡
%% point: 摩擦力
%% method: 无
%% score: 4
%% degree: 650
%% duplicateId: 0
%% body:
下列关于物体受静摩擦力作用的叙述中，正确的是\xzanswer{C}
\fourchoices[answer=C]
{静摩擦力的方向一定与物体的运动方向相反}
{静摩擦力的方向不可能与物体的运动方向相同}
{静摩擦力的方向可能与物体的运动方向垂直}
{静止物体所受静摩擦力一定为零}
%% answer: C
%% memo:
\memoanswer{
本题原卷及材料中未提供详解，待补充。
}

\item
%% number: 5
%% paperName: 1993年上海高考第 5 题 4 分
%% typeId: 1
%% type: 单选题
%% chapter: 力物体的平衡
%% point: 物体系平衡
%% method: 无
%% score: 4
%% degree: 450
%% duplicateId: 0
%% body:
如图所示，两根细线挂着两个质量相同的小球 A、B，上、下两根细线中的拉力分别是 $T_{A}$、$T_{B}$。现在使 A、B 带同种电荷，此时，上、下两根细线中的拉力分别是 $T_{A}'$、$T_{B}'$，则\xzanswer{A}
\onepicture{\includesvg{figs/05}}
\fourchoices[answer=A]
{$T_{A}' = T_{A}$，$T_{B}' > T_{B}$}
{$T_{A}' = T_{A}$，$T_{B}' < T_{B}$}
{$T_{A}' < T_{A}$，$T_{B}' > T_{B}$}
{$T_{A}' > T_{A}$，$T_{B}' < T_{B}$}
%% answer: A
%% memo:
\memoanswer{
本题原卷及材料中未提供详解，待补充。
}

\item
%% number: 6
%% paperName: 1993年上海高考第 6 题 4 分
%% typeId: 1
%% type: 单选题
%% chapter: 电场
%% point: 电荷守恒定律
%% method: 无
%% score: 4
%% degree: 550
%% duplicateId: 0
%% body:
如图所示，半径相同的两个金属小球 A、B 带有电量相等的电荷，相隔一定距离，两球之间的相互吸引力的大小是 $F$。今让第三个半径相同不带电的金属小球先后与 A、B 两球接触后移开。这时，A、B 两球之间的相互作用力的大小是\xzanswer{A}
\onepicture{\includesvg{figs/06}}
\fourchoices[answer=A]
{$\frac{1}{8}F$}
{$\frac{1}{4}F$}
{$\frac{3}{8}F$}
{$\frac{3}{4}F$}
%% answer: A
%% memo:
\memoanswer{
本题原卷及材料中未提供详解，待补充。
}

\item
%% number: 7
%% paperName: 1993年上海高考第 7 题 4 分
%% typeId: 1
%% type: 单选题
%% chapter: 磁场
%% point: 带电粒子在复合场中的运动
%% method: 无
%% score: 4
%% degree: 450
%% duplicateId: 0
%% body:
如图所示，一束质量、速度和电量不同的正离子垂直地射入匀强磁场和匀强电场正交的区域里，结果发现有些离子保持原来的运动方向，未发生任何偏转。如果让这些不偏转离子进入另一匀强磁场中，发现这些离子又分裂成几束，对这些进入后一磁场的离子，可得出结论\xzanswer{D}
\onepicture{\includesvg{figs/07}}
\fourchoices[answer=D]
{它们的动能一定各不相同}
{它们的电量一定各不相同}
{它们的质量一定各不相同}
{它们的电量和质量之比一定各不相同}
%% answer: D
%% memo:
\memoanswer{
本题原卷及材料中未提供详解，待补充。
}

\item
%% number: 8
%% paperName: 1993年上海高考第 8 题 4 分
%% typeId: 1
%% type: 单选题
%% chapter: 气体性质
%% point: 气体图象
%% method: 无
%% score: 4
%% degree: 450
%% duplicateId: 0
%% body:
如图所示，一定质量的理想气体，由状态 a 沿直线变化到状态 b，在此过程中，气体分子平均速率的变化情况\xzanswer{D}
\onepicture{\includesvg{figs/08}}
\fourchoices[answer=D]
{不断增大}
{不断减小}
{先减小后增大}
{先增大后减小}
%% answer: D
%% memo:
\memoanswer{
本题原卷及材料中未提供详解，待补充。
}

\item
%% number: 9
%% paperName: 1993年上海高考第 9 题 5 分
%% typeId: 2
%% type: 多选题
%% chapter: 气体性质
%% point: 分子力
%% method: 无
%% score: 5
%% degree: 650
%% duplicateId: 0
%% body:
下面的叙述正确的是\xzanswer{ABD}
\fourchoices[answer=ABD]
{分子之间既有引力作用，又有斥力作用}
{当分子之间距离增大时，分子间的引力和斥力都减小}
{气体分子平均动能越大，其压强一定越大}
{温度相同时，分子质量不同的两种气体，其分子平均动能一定相同}
%% answer: ABD
%% memo:
\memoanswer{
本题原卷及材料中未提供详解，待补充。
}

\item
%% number: 10
%% paperName: 1993年上海高考第 10 题 5 分
%% typeId: 2
%% type: 多选题
%% chapter: 原子和原子核
%% point: 玻尔模型
%% method: 无
%% score: 5
%% degree: 650
%% duplicateId: 0
%% body:
氢原子的核外电子由一个轨道跃迁到另一轨道时，可能发生的情况有\xzanswer{BC}
\fourchoices[answer=BC]
{放出光子，电子动能减少，原子势能增加}
{放出光子，电子动能增加，原子势能减少}
{吸收光子，电子动能减少，原子势能增加}
{吸收光子，电子动能增加，原子势能减少}
%% answer: BC
%% memo:
\memoanswer{
本题原卷及材料中未提供详解，待补充。
}

\item
%% number: 11
%% paperName: 1993年上海高考第 11 题 5 分
%% typeId: 2
%% type: 多选题
%% chapter: 光学
%% point: 透镜
%% method: 无
%% score: 5
%% degree: 450
%% duplicateId: 0
%% body:
两束与主轴距离相等的单色可见光 A 和 B，平行于主轴射向凸透镜，经透镜折射后相交于 P 点，如图所示，由此可得出\xzanswer{AD}
\onepicture{\includesvg{figs/11}}
\fourchoices[answer=AD]
{A 光在透镜玻璃中的速度比 B 光小}
{透镜玻璃对 A 光的折射率比 B 光小}
{B 光的波长比 A 光短}
{B 光的光子能量比 A 光小}
%% answer: AD
%% memo:
\memoanswer{
本题原卷及材料中未提供详解，待补充。
}

\item
%% number: 12
%% paperName: 1993年上海高考第 12 题 5 分
%% typeId: 2
%% type: 多选题
%% chapter: 交流电
%% point: LC振荡回路
%% method: 无
%% score: 5
%% degree: 450
%% duplicateId: 0
%% body:
如图表示 LC 振荡电路某时刻的情况，以下说法正确的是\xzanswer{BCD}
\onepicture{\includesvg{figs/12}}
\fourchoices[answer=BCD]
{电容器正在充电}
{电感线圈中的磁场能正在增加}
{电感线圈中的电流正在增加}
{此时刻自感电动势正在阻碍电流增大}
%% answer: BCD
%% memo:
\memoanswer{
本题原卷及材料中未提供详解，待补充。
}

\item
%% number: 13
%% paperName: 1993年上海高考第 13 题 5 分
%% typeId: 2
%% type: 多选题
%% chapter: 电路
%% point: 动态分析
%% method: 无
%% score: 5
%% degree: 350
%% duplicateId: 0
%% body:
如图所示电路，开关 K 原来是闭合的，当 $R_{1}$、$R_{2}$ 的滑片刚好处于各自的中点位置时，悬在空气平板电容器 $C$ 两水平极板间的带电尘埃 P 恰好处于静止状态。要使尘埃 P 加速向上运动的方法是\xzanswer{CD}
\onepicture{\includesvg{figs/13}}
\fourchoices[answer=CD]
{把 $R_{1}$ 的滑片向上移动}
{把 $R_{2}$ 的滑片向上移动}
{把 $R_{2}$ 的滑片向下移动}
{把开关 K 断开}
%% answer: CD
%% memo:
\memoanswer{
本题原卷及材料中未提供详解，待补充。
}

\item
%% number: 14
%% paperName: 1993年上海高考第 14 题 4 分
%% typeId: 3
%% type: 填空题
%% chapter: 光学
%% point: 光子说
%% method: 无
%% score: 4
%% degree: 850
%% duplicateId: 0
%% body:
17 世纪，物理学家 \tkanswer[1.6cm]{惠更斯} 提出了光的波动说。为了解释 \tkanswer[2.2cm]{光电效应} 实验，爱因斯坦提出了光子说。
%% answer: 
\jdanswer{
惠更斯，光电效应
}
%% memo:
\memoanswer{
本题原卷及材料中未提供详解，待补充。
}

\item
%% number: 15
%% paperName: 1993年上海高考第 15 题 4 分
%% typeId: 3
%% type: 填空题
%% chapter: 功和能
%% point: 机械能守恒定律
%% method: 无
%% score: 4
%% degree: 550
%% duplicateId: 0
%% body:
枪竖直向上以初速 $v_{0}$ 发射子弹，忽略空气阻力，当子弹离枪口距离为 \tkanswer[2.4cm]{$\frac{v_{0}^{2}}{3g}$} 时，子弹的动能是其重力势能的一半。
%% answer: 
\jdanswer{
$\frac{v_{0}^{2}}{3g}$
}
%% memo:
\memoanswer{
本题原卷及材料中未提供详解，待补充。
}

\item
%% number: 16
%% paperName: 1993年上海高考第 16 题 4 分
%% typeId: 3
%% type: 填空题
%% chapter: 运动和力
%% point: 瞬时加速度
%% method: 无
%% score: 4
%% degree: 550
%% duplicateId: 0
%% body:
如图所示，木块 A 和 B 用一轻弹簧相连，竖直放在木块 C 上，三者静置于地面，它们的质量之比是 $1:2:3$。设所有接触面都光滑，当沿水平方向迅速抽出木块 C 的瞬时，A 和 B 的加速度分别是 $a_{A} =$ \tkanswer[1.4cm]{$0$}，$a_{B} =$ \tkanswer[1.8cm]{$1.5g$}。
\onepicture[align=right]{\includesvg{figs/16}}
%% answer: 
\jdanswer{
$0$，$1.5g$
}
%% memo:
\memoanswer{
本题原卷及材料中未提供详解，待补充。
}

\item
%% number: 17
%% paperName: 1993年上海高考第 17 题 4 分
%% typeId: 3
%% type: 填空题
%% chapter: 曲线运动
%% point: 平抛运动
%% method: 无
%% score: 4
%% degree: 350
%% duplicateId: 0
%% body:
如图所示，高为 $h$ 的车厢在平直轨道上匀减速向右行驶，加速度大小为 $a$，车厢顶部 A 点处有油滴滴落到车厢地板上，车厢地板上的 O 点位于 A 点的正下方，则油滴落地点必在 O 点的 \tkanswer[1.2cm]{右}（填左、右）方，离 O 点距离为 \tkanswer[2cm]{$\frac{ah}{g}$}。
\onepicture[align=right]{\includesvg{figs/17}}
%% answer: 
\jdanswer{
右，$\frac{ah}{g}$
}
%% memo:
\memoanswer{
本题原卷及材料中未提供详解，待补充。
}

\item
%% number: 18
%% paperName: 1993年上海高考第 18 题 4 分
%% typeId: 3
%% type: 填空题
%% chapter: 机械波
%% point: 波的图象
%% method: 无
%% score: 4
%% degree: 450
%% duplicateId: 0
%% body:
如图所示，一列横波沿 $x$ 轴正方向传播，当位于 $x_{1} = 2\Ucm$ 处的 A 质点处在 $x$ 轴的下方最大位移处时，位于 $x_{2} = 5\Ucm$ 处的 B 质点恰好在平衡位置，而且运动方向向上。已知 A、B 两质点在 $x$ 轴的距离小于波长，则这列波的波长为 \tkanswer[1.8cm]{$4\Ucm$}，在图中画出 $x = 0$ 到 $x = 8\Ucm$ 区间内该时刻波的图象。
\onepicture[align=right]{\includesvg{figs/18}}
%% answer: 
\jdanswer{
\begin{enumerate}
	\item
	$4\Ucm$

	\item
	\onepicture[w=7cm]{figs/18_an.png}

\end{enumerate}
}
%% memo:
\memoanswer{
本题原卷及材料中未提供详解，待补充。
}

\item
%% number: 19
%% paperName: 1993年上海高考第 19 题 4 分
%% typeId: 3
%% type: 填空题
%% chapter: 功和能
%% point: 动能定理
%% method: 无
%% score: 4
%% degree: 450
%% duplicateId: 0
%% body:
如图所示，物体质量为 $m$，由静止开始从 A 点沿斜面从 $h_{1}$ 高处下滑到地面，随后又沿另一斜面上滑到 $h_{2}$ 高处的 B 点停止，若在 B 点给物体一瞬时冲量，使物体从 B 点沿原路返回到 A 点，需给物体的最小冲量值是 \tkanswer[2.6cm]{$2m\sqrt{g(h_{1}-h_{2})}$}。
\onepicture[align=right]{\includesvg{figs/19}}
%% answer: 
\jdanswer{
$2m\sqrt{g(h_{1}-h_{2})}$
}
%% memo:
\memoanswer{
本题原卷及材料中未提供详解，待补充。
}

\item
%% number: 20
%% paperName: 1993年上海高考第 20 题 4 分
%% typeId: 3
%% type: 填空题
%% chapter: 磁场
%% point: 带电粒子在复合场中的运动
%% method: 无
%% score: 4
%% degree: 450
%% duplicateId: 0
%% body:
如图所示，带电液滴从 $h$ 高处自由落下，进入一个匀强电场与匀强磁场互相垂直的区域，磁场方向垂直纸面，电场强度为 $E$，磁感应强度为 $B$，已知液滴在此区域中作匀速圆周运动，则圆周的半径 $R =$ \tkanswer[2.8cm]{$\frac{E}{B}\sqrt{\frac{2h}{g}}$}。
\onepicture[align=right]{\includesvg{figs/20}}
%% answer: 
\jdanswer{
$\frac{E}{B}\sqrt{\frac{2h}{g}}$
}
%% memo:
\memoanswer{
本题原卷及材料中未提供详解，待补充。
}

\item
%% number: 20
%% paperName: 1993年上海高考第 20 题 4 分
%% typeId: 3
%% type: 填空题
%% chapter: 电路
%% point: 电容
%% method: 无
%% score: 4
%% degree: 450
%% duplicateId: 0
%% body:
右图中 $E = 10\UV$，$C_{1} = C_{2} = 30\UmuF$，$R_{1} = 4.0\UO$，$R_{2} = 6.0\UO$，电池内阻可忽略。先闭合开关 K，待电路稳定后，再将 K 断开，则断开后流过电阻 $R_{1}$ 的电量为 \tkanswer[2.6cm]{$4.2\times10^{-4}\UC$}。
\onepicture[align=right]{figs/21.png}
%% answer: 
\jdanswer{
$4.2\times10^{-4}\UC$
}
%% memo:
\memoanswer{
本题原卷及材料中未提供详解，待补充。
}

\item
%% number: 22
%% paperName: 1993年上海高考第 22 题 4 分
%% typeId: 4
%% type: 实验题
%% chapter: 固体液体的性质
%% point: 表面张力
%% method: 无
%% score: 4
%% degree: 650
%% duplicateId: 0
%% body:
右图金属框上阴影部分表示肥皂膜，它被棉线分割成 a、b 两部分，若将肥皂膜的 a 部分用热针刺破，棉线的形状是下图中的哪一个？\xzanswer{D}
\onepicture[align=right]{\includesvg{figs/22a}}
\fourchoices[answer=D, ispicture=true, h=3.2cm]
{figs/22b1.png}
{figs/22b2.png}
{figs/22b3.png}
{figs/22b4.png}
%% answer: D
\jdanswer{
D
}
%% memo:
\memoanswer{
本题原卷及材料中未提供详解，待补充。
}

\item
%% number: 23
%% paperName: 1993年上海高考第 23 题 5 分
%% typeId: 4
%% type: 实验题
%% chapter: 运动和力
%% point: 验证牛顿第二定律
%% method: 无
%% score: 5
%% degree: 550
%% duplicateId: 0
%% body:
在验证牛顿第二定律关于作用力一定时，加速度与质量成反比的实验中，以下做法错误的是\xzanswer{ACD}
\fourchoices[answer=ACD]
{平衡摩擦力时，应将装沙的小桶用细绳通过定滑轮系在小车上}
{每次改变小车的质量时，不需要重新平衡摩擦力}
{实验时，先放开小车，再接通打点计时器电源}
{小车运动的加速度可从天平测出装沙小桶和沙的质量（$M'$ 与 $m'$）以及小车质量 $M$，直接用公式 $a = \frac{M' + m'}{M}g$ 求出。$[(M' + m') \ll M]$}
%% answer: ACD
\jdanswer{
ACD
}
%% memo:
\memoanswer{
本题原卷及材料中未提供详解，待补充。
}

\item
%% number: 24
%% paperName: 1993年上海高考第 24 题 5 分
%% typeId: 4
%% type: 实验题
%% chapter: 气体性质
%% point: 验证玻意耳定律
%% method: 无
%% score: 5
%% degree: 650
%% duplicateId: 0
%% body:
在《验证玻意耳——马略特定律》的实验中，下列操作哪些是不正确的？\xzanswer{BCD}
\fourchoices[answer=BCD]
{实验前，记下室内大气压强的数值}
{先将橡皮帽封住注射小孔，然后再将活塞推入针筒}
{在实验过程中，若橡皮帽突然脱落，应立即将它堵住}
{框架两侧挂上钩码后，立即记下空气柱的长度}
%% answer: BCD
\jdanswer{
BCD
}
%% memo:
\memoanswer{
本题原卷及材料中未提供详解，待补充。
}

\item
%% number: 25
%% paperName: 1993年上海高考第 25 题 4 分
%% typeId: 4
%% type: 实验题
%% chapter: 力物体的平衡
%% point: 力矩盘
%% method: 无
%% score: 4
%% degree: 550
%% duplicateId: 0
%% body:
演示力矩盘平衡的实验图如图，图中 O 是固定转动轴，A、B 是两个光滑的定滑轮，两组砝码跨过定滑轮悬挂，有一组砝码原悬挂在 C 点，实验时不慎将其悬线跨过 P 点处的大头针，成图中所示的情况。力矩盘处于平衡状态。盘上每个圆环间距都是 $2.0\Ucm$，每个砝码重 $0.10\UN$。则从图中可得出顺时针方向的力矩之和是 \tkanswer[2.4cm]{$0.04\UNcdotm$}。实验时，由于力矩盘偏心未经调整实际测出的弹簧秤读数偏大，则力矩盘的重心在轴的 \tkanswer[1.6cm]{右方}。（填左方或右方）
\onepicture[align=right]{\includesvg{figs/25}}
%% answer: 
\jdanswer{
$0.04\UNcdotm$，右方
}
%% memo:
\memoanswer{
本题原卷及材料中未提供详解，待补充。
}

\item
%% number: 26
%% paperName: 1993年上海高考第 26 题 6 分
%% typeId: 4
%% type: 实验题
%% chapter: 电路
%% point: 串并联电路
%% method: 无
%% score: 6
%% degree: 450
%% duplicateId: 0
%% body:
有小灯泡 $L_{1}$（$6\UV$，$3\UW$）和 $L_{2}$（$6\UV$，$1\UW$）各一只，定值电阻 $R$（$18\UO$，$5\UW$）一只，电源一只（电动势 $12\UV$，内阻不计），电键一只，导线若干。试设计一电路，使 $L_{1}$ 和 $L_{2}$ 都能正常发光。
\begin{steps}
	\item
	将设计的电路图画在虚线框内，电路图中要标明各元件的代号。
	\item
	按设计的电路图，用线（表示导线）完成下图所示仪器的连线。
\end{steps}
\onepicture[align=right]{\includesvg{figs/26}}
%% answer: 
\jdanswer{
\begin{enumerate}
	\item
	设计的电路图如图：
	\onepicture[w=5cm]{figs/26_an1.png}

	\item
	把实物连接成的电路如图：
	\onepicture[w=8cm]{figs/26_an2.png}

\end{enumerate}
}
%% memo:
\memoanswer{
本题原卷及材料中未提供详解，待补充。
}

\item
%% number: 27
%% paperName: 1993年上海高考第 27 题 12 分
%% typeId: 6
%% type: 计算题
%% chapter: 电磁感应
%% point: 线圈过有界磁场
%% method: 无
%% score: 12
%% degree: 450
%% duplicateId: 0
%% body:
如图所示，Ⅰ、Ⅲ 为两匀强磁场区，Ⅰ 区域的磁场方向垂直纸面向里，Ⅲ 区域的磁场方向垂直纸面向外，磁感应强度均为 $B$，两区域中间为宽 $s$ 的无磁场区 Ⅱ。有一边长为 $L$（$L > s$），电阻为 $R$ 的正方形金属框 abcd 置于 Ⅰ 区域，ab 边与磁场边界平行，现拉着金属框以速度 $v$ 向右匀速移动。
\begin{enumerate}
	\item
	分别求出当 ab 边刚进入中央无磁场区 Ⅱ，和刚进入磁场区 Ⅲ 时，通过 ab 边的电流的大小和方向。
	\item
	把金属框从 Ⅰ 区域完全拉入 Ⅲ 区域过程中拉力所做的功。
\end{enumerate}
\onepicture[align=right]{\includesvg{figs/27}}
%% answer: 
\jdanswer{
\begin{enumerate}
	\item
	$I_{1} = \frac{BLv}{R}$，方向 b$\to$a；$I_{2} = \frac{2BLv}{R}$，方向 b$\to$a

	\item
	$W = \frac{2B^{2}L^{2}v(2L - s)}{R}$

\end{enumerate}
}
%% memo:
\memoanswer{
（1）当 ab 边刚进入中间无磁场区 Ⅱ 时的电流：

$I_{1} = \frac{BLv}{R}$，方向 b$\to$a

刚进入磁场区 Ⅲ 时

$I_{2} = \frac{2BLv}{R}$，方向 b$\to$a

（2）在 ab 边穿过宽 $s$ 的 Ⅱ 区过程中，cd 边受安培力 $F_{1} = BI_{1}l = \frac{B^{2}L^{2}v}{R}$，拉力做功 $W_{1} = F_{1}s = \frac{B^{2}L^{2}v}{R}s$；

在 ab 边进入 Ⅲ 区，$F_{1} = BI_{2}l = \frac{2B^{2}L^{2}v}{R}$，匀速拉动外力应等于 $2F_{2}$，通过距离为 $(L - s)$，拉力做功 $W_{2} = 2F_{2}(L - s) = \frac{4B^{2}L^{2}v(L - s)}{R}$；

在 cd 边通过 Ⅲ 区域过程中，只有 ab 边受安培力，且 $F_{3} = F_{1}$，距离仍为 $s$，拉力做功 $W_{3} = W_{1}$。在线框完全进入 Ⅲ 后无感应电流，不受磁场力，拉力做功为零，所以总功为：

$W = W_{1} + W_{2} + W_{3} = \frac{2B^{2}L^{2}v(2L - s)}{R}$

评分标准：本题共 12 分，（1）6 分，（2）6 分。

（1）正确得出 $I_{1}$、$I_{2}$ 各 2 分，正确表示方向各 1 分。

（2）正确得出 $W_{1}$、$W_{2}$ 各 2 分，正确得出 $W_{3}$ 得 1 分，正确得出 $W$ 得 1 分。
}

\item
%% number: 28
%% paperName: 1993年上海高考第 28 题 10 分
%% typeId: 6
%% type: 计算题
%% chapter: 光学
%% point: 透镜
%% method: 无
%% score: 10
%% degree: 550
%% duplicateId: 0
%% body:
如图所示，半径 $r = 1.0\Ucm$ 的圆盘形发光面 S 和半径相同的不透光圆盘 D 共轴放置，相距 $8.0\Ucm$，在 D 的右方 $4.0\Ucm$ 处有一光屏 P，如图所示。
\begin{enumerate}
	\item
	求屏上 D 盘本影区域的半径 $R_{1}$ 和半影区域的外半径 $R_{2}$。
	\item
	若在 S 与 D 的正中间 O 处共轴放置一焦距为 $8.0\Ucm$ 的凸透镜，求此时屏上 D 盘本影区域的半径 $R_{1}'$ 和半影区域的外半径 $R_{2}'$。（设透镜的半径远大于 $r$）
\end{enumerate}
\onepicture[align=right]{\includesvg{figs/28}}
%% answer: 
\jdanswer{
\begin{enumerate}
	\item
	$R_{1} = 1.0\Ucm$，$R_{2} = 2.0\Ucm$

	\item
	$R_{1}' = 0.67\Ucm$，$R_{2}' = 2.0\Ucm$

\end{enumerate}
}
%% memo:
\memoanswer{
（1）$R_{1} = 1.0\Ucm$，$R_{2} = 2.0\Ucm$。如右图。

\onepicture[align=right, w=8cm]{figs/28_an.png}

（2）S 经透镜成虚像在 $S'$ 处（如图）。

以 $u = 4$，$f = 8$ 代入公式

$\frac{1}{u} + \frac{1}{v} = \frac{1}{f}$

得 $v = -8\Ucm$

$S'$ 得半径 $r' = \frac{\left| v \right|}{u}\cdot r = 2\Ucm$

发光面 S 经透镜折射后，其光线如从其像 $S'$ 发出，故由几何关系：

$\frac{R_{1}'}{r} = \frac{8}{12}$，$R_{1}' = \frac{8}{12}r = 0.67\Ucm$。

$\frac{R_{2}'}{r} = \frac{8}{4}$，$R_{2}' = \frac{8}{4}r = 2.0\Ucm$。

评分标准：本题共 10 分。（1）4 分，（2）6 分。

（1）正确求出 $R_{1}$，$R_{2}$ 各 2 分。

（2）正确求出虚像位置 $v$ 和大小 $r'$ 各得 1 分。正确得出 $R_{1}'$，$R_{2}'$ 各 2 分。
}

\item
%% number: 29
%% paperName: 1993年上海高考第 29 题 15 分
%% typeId: 6
%% type: 计算题
%% chapter: 功和能
%% point: 子弹打木块模型
%% method: 分情况讨论
%% score: 15
%% degree: 350
%% duplicateId: 0
%% body:
如图所示，质量为 $2m$、长为 $l$ 的木块置于光滑水平台面上，质量为 $m$ 的子弹以初速 $v_{0}$ 水平向右射向木块，穿出木块时的速度为 $\frac{v_{0}}{2}$，设木块对子弹的阻力恒定。
\begin{enumerate}
	\item
	求子弹穿越木块的过程中木块滑行的距离 $L_{1}$；
	\item
	若木块固定在传送带上，使木块随传送带始终以恒定速度 $u$ 水平向右运动，子弹仍以速度 $v_{0}$ 向右射向木块（$u < v_{0}$），求子弹最终速度 $v$；
	\item
	求在（2）情况下，子弹在木块中行进的过程中，木块移动的距离 $L_{2}$。
\end{enumerate}
\onepicture[align=right]{\includesvg{figs/29}}
%% answer: 
\jdanswer{
\begin{enumerate}
	\item
	$L_{1} = \frac{1}{5}l$

	\item
	$v = u + \sqrt{(v_{0} - u)^{2} - \frac{5}{8}v_{0}^{2}}$，当 $u \leq v_{0}\left(1 - \frac{\sqrt{10}}{4}\right)$；$v = u$，当 $v_{0} > u > v_{0}\left(1 - \frac{\sqrt{10}}{4}\right)$

	\item
	当 $u \leq v_{0}\left(1 - \frac{\sqrt{10}}{4}\right)$，$L_{2} = \frac{16ul}{5v_{0}^{2}}\left[v_{0} - u - \sqrt{(v_{0} - u)^{2} - \frac{5}{8}v_{0}^{2}}\right]$；当 $v_{0} > u > v_{0}\left(1 - \frac{\sqrt{10}}{4}\right)$，$L_{2} = \frac{16ul}{5v_{0}^{2}}(v_{0} - u)$

\end{enumerate}
}
%% memo:
\memoanswer{
（1）子弹与木块组成的系统动量守恒，则有

$mv_{0} = 2mV + \frac{mv_{0}}{2}$ ①

得 $V = \frac{v_{0}}{4}$

设木块对子弹阻力为 $f$，对子弹和木块分别应用动能定理，有

$-f(L_{1} + l) = \frac{1}{2}m\left(\frac{v_{0}}{2}\right)^{2} - \frac{1}{2}mv_{0}^{2}$ ②

$fL_{1} = \frac{1}{2}(2m)V^{2}$ ③

两式相加得 $fl = \frac{5}{16}mv_{0}^{2}$ ④

由③、④消去 $f$，以 $V = \frac{v_{0}}{4}$ 代入，得 $L_{1} = \frac{1}{5}l$

（2）对子弹应用动能定理：

$-f(l + ut) = \frac{1}{2}mv^{2} - \frac{1}{2}mv_{0}^{2}$ ⑤

对子弹应用动量定理：

$-ft = mv - mv_{0}$ ⑥

由⑥ $fut = mu(v_{0} - v)$

代入⑤ $-fl - mu(v_{0} - v) = \frac{1}{2}mv^{2} - \frac{1}{2}mv_{0}^{2}$

以④代入上式，即可解得 $v = u + \sqrt{(v_{0} - u)^{2} - \frac{5}{8}v_{0}^{2}}$ ⑦

此解只在 $(v_{0} - u)^{2} \geq \frac{5}{8}v_{0}^{2}$ 即 $u \leq v_{0}\left(1 - \frac{\sqrt{10}}{4}\right)$ 时成立。

当 $u > v_{0}\left(1 - \frac{\sqrt{10}}{4}\right)$ 时，子弹不再从木块射出，子弹的最终速度为 $u$，故结果为：

$v = u + \sqrt{(v_{0} - u)^{2} - \frac{5}{8}v_{0}^{2}}$，当 $u \leq v_{0}\left(1 - \frac{\sqrt{10}}{4}\right)$；

$v = u$，当 $v_{0} > u > v_{0}\left(1 - \frac{\sqrt{10}}{4}\right)$

（3）当 $u \leq v_{0}\left(1 - \frac{\sqrt{10}}{4}\right)$ 时，子弹在木块中行进时间 $t$ 可由⑥式得到：

$t = \frac{m}{f}(v_{0} - v)$

以④、⑦代入上式得 $t = \frac{16l}{5v_{0}^{2}}\left[v_{0} - u - \sqrt{(v_{0} - u)^{2} - \frac{5}{8}v_{0}^{2}}\right]$

由此得 $L_{2} = ut = \frac{16ul}{5v_{0}^{2}}\left[v_{0} - u - \sqrt{(v_{0} - u)^{2} - \frac{5}{8}v_{0}^{2}}\right]$

当 $v_{0} > u > v_{0}\left(1 - \frac{\sqrt{10}}{4}\right)$ 时，子弹在木块中行进时间

$t' = \frac{m}{f}(v_{0} - u) = \frac{16l}{5v_{0}^{2}}(v_{0} - u)$

$L_{2}' = ut' = \frac{16ul}{5v_{0}^{2}}(v_{0} - u)$

评分标准：本题共 15 分。（1）5 分，（2）6 分，（3）4 分。

（1）正确列出①、②、③式各 1 分。正确得出 $L_{1}$ 结果得 2 分。

（2）正确列出⑤、⑥式各 1 分，正确得出 $v$ 的两个结果各 1 分，正确写出各结果适用条件（$u$ 的取值）各 1 分。

（3）正确得出 $L_{2}$、$L_{2}'$ 各 1 分，正确写出各结果使用条件各 1 分。
}

\end{enumerate}

\end{document}
